% FILE: 3-poster/parts/colors.tex
% ==================================================================================================
% Poster palette. THIS FILE IS THE ONLY PLACE A COLOUR MAY BE DEFINED. No call site may write an
% \definecolor, an HTML value, or a hand-mixed shade; a call site names a token and nothing else.
%
% THREE INDEPENDENT PALETTES. Each token belongs to exactly ONE of them, so any palette can be
% retuned without shifting the other two:
%   1. Neutral foundation — the achromatic substrate: background, surfaces, borders, text levels, and
%                           the neutral structure of plots (axes, reference rules). Shared by chrome
%                           AND figures, but never an accent and never a concept.
%   2. Design chrome      — turquoise (+ a reserved complementary coral): headings, boxes, navigation,
%                           QR tags, badges, and the two kinds of POINTER AWAY from the poster
%                           (\litcite, \seebottom) -- one colour, because they are one kind of mark.
%   3. Semantic grammar   — one concept -> one colour, used ONLY inside figures.
%
% >>> THE ONE LICENSED EXCEPTION TO THE CHROME RULE <<<
% The chrome turquoise used to be "NEVER used to encode a concept inside a figure". It now has exactly
% one concept: `netchrome' marks THE POSTER'S OWN ACCOUNT (the recurrent network model) in the
% model-comparison band. That is deliberate -- see "TWO MODELS" below -- and it is the only exception.
% The rule holds everywhere else.
%
% Values are sRGB approximations of an OKLCH specification. The perceptual ladder (relative lightness
% ordering within a hue) is what matters, not exact colorimetry.

% ----- 1. Neutral foundation (achromatic substrate; independent of both accent palettes) -----------
\definecolor{posterbg}{HTML}{F1F3F3}      % L=0.970  background
\definecolor{surface}{HTML}{FBFCFC}       % L=0.995  elevated panels (shadow separates from bg)
\definecolor{bordersubtle}{HTML}{DBDFDF}  % L=0.900  decorative outlines, dividers, table lines
\definecolor{borderfunc}{HTML}{8C9698}    % L<=0.640 focus rings, emphasis boundaries, plot axes
\definecolor{texttertiary}{HTML}{6E797B}  % L=0.530  captions, minor metadata, faint hints
\definecolor{textsecondary}{HTML}{586366} % L=0.450  subtitles, supporting text
\definecolor{textprimary}{HTML}{1B2224}   % L=0.150  titles, labels, body, key values

% ----- 2. Design-chrome accent (headings, boxes, navigation) ---------------------------------------
\definecolor{accentturq}{HTML}{12A89E}    % primary design accent (H ~ 187); also \litcite + \seebottom
\definecolor{accentturqdk}{HTML}{0C7A72}  % darker turquoise: sub-headers, contact line, badges
\definecolor{accentcoral}{HTML}{F26B5B}   % complementary DESIGN accent -- RESERVED (chrome only)

% ==================================================================================================
% 3. SEMANTIC GRAMMAR — one concept, one colour. Used ONLY inside figures.
% ==================================================================================================
%
% ----- 3a. THE TWO TRANSFORMATIONS: a blue pole and a magenta pole ---------------------------------
% WHY THEY ARE NOT ORANGE AND GREEN ANY MORE, which is the defect this palette exists to correct.
% The two components are red and green (compred / compgreen). The two transformations used to be
% orange and green -- and MEASURED IN OKLab the amplification colour sat at dE 0.107 from component 2,
% below the ~0.15 a reader needs to tell two categories apart. The projection sat at 0.156 from
% component 1, i.e. exactly on the threshold. The reader was being asked to hold "green = the second
% component" and "green = the recurrent amplification" at once, in figures where both appear.
% The components name WHAT IS COMBINED; the transformations name HOW THE NETWORK ACTS. Two concepts of
% different kinds must occupy different hue REGIONS, not adjacent ones. Blue/magenta vacates the
% red/green region entirely. Blue reads as exogenous, informational, directional; magenta as
% endogenous, energetic, self-reinforcing -- which is the model: the input arrives from outside, the
% recurrence feeds back on itself. (H 258 and H 330 sit 72 deg apart: a strong contrast, NOT an
% opponent complement. The pairing carries its meaning through warm/cool and hue distance.)
%
% ONE LIGHTNESS LADDER, SHARED BY BOTH POLES. Each transformation is a constant-hue tint of its own
% SOURCE concept, so "projection" is visibly the input's operation and "amplification" is visibly the
% recurrence's operation.
%
%   level   role                                   blue (H 258)   magenta (H 330)
%   deep    card chrome / banners (white text)     #0F4EA3        #7D2A79
%   core    the SOURCE concept (input I; conn. J)  #2368C9        #9E4098
%   bright  the TRANSFORMATION (xi_k ; A_k)        #3483F3        #C154BA
%   light   a second pattern within a pole         #74A6EF        #CF88C8
%   pale    fills and washes (needs a stroke)      #D7E6FC        #F3DDF0
%
% >>> THE LADDER LIMIT, AND IT BINDS <<<
% ADJACENT RUNGS OF ONE HUE SIT AT dE ~ 0.09 -- BELOW the categorical threshold. Only NON-ADJACENT
% rungs are reliably distinct (deep vs bright: dE 0.186). Two rules follow:
%   (1) Within ONE figure, use at most THREE levels of one hue and SKIP A RUNG between them
%       (deep, bright, pale). NEVER place two adjacent rungs in the same figure.
%   (2) The two pale fills sit at dE 0.041 from EACH OTHER, so PALE CANNOT ENCODE THE POLE. A pale
%       fill must ALWAYS be paired with a saturated stroke of its own family.
% If a figure needs a fourth distinction, move it to a NON-COLOUR channel. The index of a pattern
% (first or second) is an INDEX, not a concept: it belongs to a subscript or a line style, never to a
% fourth tint. See figures/approach-steps.tex, which is the figure this rule was written for.
%
% --- blue pole: the external input and what it does (the feedforward projection) ---
\definecolor{inputbluedk}{HTML}{0F4EA3}   % deep   -- input pattern 1 (I^(1)); also the projection card chrome
\definecolor{inputblue}{HTML}{2368C9}     % core   -- the external input itself (I, u)
\definecolor{projblue}{HTML}{3483F3}      % bright -- THE FEEDFORWARD PROJECTION (xi_k, Sigma^(nI))
\definecolor{inputbluelt}{HTML}{74A6EF}   % light  -- input pattern 2 (I^(2))
\definecolor{inpoolblue}{HTML}{D7E6FC}    % pale   -- input pooling Sigma^(II)  [FILL ONLY: needs a stroke]
% --- magenta pole: the recurrence and what it does (the recurrent amplification) ---
\definecolor{recurmagdk}{HTML}{7D2A79}    % deep   -- recurrent pattern 1 (m_1); the gain; the amp. card chrome
\definecolor{recurmag}{HTML}{9E4098}      % core   -- the recurrent connectivity itself (J, Sigma^(nm))
\definecolor{ampmag}{HTML}{C154BA}        % bright -- THE RECURRENT AMPLIFICATION (A_k)
\definecolor{recurmaglt}{HTML}{CF88C8}    % light  -- recurrent pattern 2 (m_2)
\definecolor{recpoolmag}{HTML}{F3DDF0}    % pale   -- recurrent pooling Sigma^(mm)  [FILL ONLY: needs a stroke]
% THE GAIN IS MAGENTA, and this is a claim, not a convenience. g(Delta) parameterises A_k(g): it is
% part of the ACTUATOR. Left in its old blue (#2E6FC4 -- which is now, exactly, the input hue) it would
% read as an exogenous quantity, which is the opposite of what it is. It takes the DEEP magenta so it
% stays distinct from the BRIGHT magenta of A_k in figures where both appear (amplification-curves).
\colorlet{gainmag}{recurmagdk}
% THE POOLING PAIR IS INVERTED RELATIVE TO THE OLD PALETTE. The INPUT pooling used to be warm and the
% RECURRENT pooling cold -- the exact opposite of the metaphor. Sigma^(II) pools the input: blue.
% Sigma^(mm) pools the recurrence: magenta.

% ----- 3b. The components, the mixture, and the geometry -------------------------------------------
% Untouched by the migration, and the reason for it: the red/green region is now theirs alone.
\definecolor{compred}{HTML}{D6453C}       % component response 1 (kappa_1)          H ~ 28
\definecolor{compgreen}{HTML}{4F9A3F}     % component response 2 (kappa_2)          H ~ 140
\definecolor{mixamber}{HTML}{F2A30C}      % THE RESPONSE: mixture response kappa_mix, network state x,
                                          % reduced state kappa -- ONE concept, ONE token.  H ~ 74
\definecolor{spangray}{HTML}{B7BFC1}      % components plane / proportional reference
\definecolor{emergentgray}{HTML}{566063}  % out-of-plane annotations (bias arc, in-plane foot marker)
\definecolor{emergentbrown}{HTML}{8A5524} % orthogonal component (eta_perp) -- solid brown, all figures
% THE RESPONSE LEFT PURPLE. `collectivepurple' (H 307) fell BETWEEN the two new poles, so the response
% vector would have read as a blend of the projection and the amplification rather than as their
% output. It is unified onto mixamber -- which already carried the mixture response in the geometry
% figures, so the merge also retires a duplicate encoding of one concept and vacates the hue slot.

% ----- 3c. THE TWO MODELS: a chromatic/achromatic split, not a third hue pair ----------------------
% NO HUE REMAINS ON THE CIRCLE, and that is measurable rather than a matter of taste. After the
% migration the chromatic tokens sit at H = 28 (component 1), 60 (eta_perp), 74 (the response),
% 140 (component 2), 187 (chrome), 258 (feedforward), 330 (amplification). The widest gap left spans
% 72 deg, leaving +/-36 deg of clearance -- under the ~45 deg a categorical hue needs at this chroma.
% A violet forced into that gap lands at dE 0.13 from BOTH transformations: it would read as a blend
% of the projection and the amplification, which is the exact defect that retired collectivepurple.
% The neutral lane holds only ONE model: an even-handed achromatic PAIR fails too, because its light
% member sits dE 0.034 from borderfunc (the axis colour) and reaches only 2.4:1 on the background,
% under the 3:1 minimum for a graphical mark.
% So the two models split on the CHROMATIC/ACHROMATIC axis, reinforced by LINE STYLE:
\definecolor{netchrome}{HTML}{00948B}     % THE RECURRENT NETWORK MODEL (this work): solid, filled
                                          % markers. It carries the poster's OWN identity colour --
                                          % the licensed exception to the chrome rule at the top of
                                          % this file. Model-comparison band only.
\definecolor{dnslate}{HTML}{474E51}       % DIVISIVE NORMALIZATION (the rival): dashed, hollow markers.
                                          % A mechanism earns a semantic hue; a descriptive fit does
                                          % not. The two differ at dE 0.206 and spend no new hue.
% This retires three colliding tokens at once: `dncoral' (dE 0.096 from compred -- and byte-identical
% to accentcoral, so a reader met red-ish marks meaning two unrelated things), `recurjade' (dE 0.079
% from the chrome turquoise) and `netturq' (dE 0.034 -- effectively the chrome already).
% KNOWN, AND RECORDED SO A LATER EDIT DOES NOT REINTRODUCE IT: dnslate falls in the same neighbourhood
% as emergentgray, which annotates the geometry sketches. THE TWO NEVER APPEAR IN ONE FIGURE. If either
% token moves, re-check that.

% ----- convenience aliases -------------------------------------------------------------------------
\colorlet{cardborder}{bordersubtle}
\colorlet{cardshadow}{black!38}   % grey drop shadow
\colorlet{cardshadowdk}{black!52} % deeper shadow for the emphasised question / results boxes
% CARD CHROME for the two transformation cards: the DEEP rung of each pole (dark enough to carry white
% text). The vivid `projblue' / `ampmag' stay for the sketches and the thin labels.
\colorlet{projbanner}{inputbluedk}   % feedforward card: ring, glow, disc, classifier pill
\colorlet{ampbanner}{recurmagdk}     % amplification card: ring, glow, disc, classifier pill
% The transformations, named by their ROLE rather than by their rung, for call sites that mean
% "the projection" / "the amplification" rather than "this particular blue".
\colorlet{projstage}{projblue}
\colorlet{ampstage}{ampmag}
